\documentclass[11pt]{article}
\usepackage{fontspec,amsmath,amssymb,amsthm,hyperref}
\usepackage[margin=1in]{geometry}
\setlength{\emergencystretch}{3em}
\title{Iterate entropy contraction and obtain the bias constant 8}
\author{LangevinLean mathematical exposition}
\date{}
\begin{document}\maketitle
\noindent Source module: \texttt{UnadjustedLangevin/LSI/DiscreteEntropy.lean}. Compile this proof document with XeLaTeX or LuaLaTeX. The accompanying HTML pairs the mathematical argument with the exact Lean statements and records the scope of the formalization.


Starting from \(H_{k+1}\le e^{-\alpha h}H_k+6h^2dL^2\), induction retains the exact geometric sum. The small-step denominator estimate \(1-e^{-\alpha h}\ge\frac34\alpha h\) gives \(H_k\le e^{-\alpha hk}H_0+8hdL^2/\alpha\).


\section{Keep the exact geometric sum before simplifying it}


Set \(q=e^{-\alpha h}\in(0,1)\), \(b=6h^2dL^2\). The affine recurrence theorem sums \(H_{k+1}\le qH_k+b\) into \(H_k\le q^kH_0+b(1-q^k)/(1-q)\). Rewrite \(q^k=e^{-\alpha hk}\) using the exponential natural-multiple identity. No sign assumption on \(H_0\) is needed.



\noindent Lean declarations: \texttt{discrete\_entropy\_recurrence\_exact}.

\section{three quarters mul le one sub exp neg}


The tangent bound \(e^x\ge1+x\), multiplied by \(e^{-x}\), gives \(1-e^{-x}\ge x/(1+x)\). For \(0\le x\le1/4\), elementary positive-denominator algebra shows \(x/(1+x)\ge3x/4\). This lower bound controls the geometric-series denominator in the bias floor.



\noindent Lean declarations: \texttt{three\_quarters\_mul\_le\_one\_sub\_exp\_neg}.

\section{discrete entropy recurrence bound}


The nonnegative forcing gives \(H_k\le e^{-\alpha hk}H_0+6h^2dL^2/(1-e^{-\alpha h})\). For \(\alpha h\le1/4\), the denominator is at least \(3\alpha h/4>0\). Divide and cancel to obtain the bias floor \(8hdL^2/\alpha\). Lean proves positivity of the denominator before multiplying across it.



\noindent Lean declarations: \texttt{discrete\_entropy\_recurrence\_bound}.

\end{document}
